\chapter{Brauer characters}
\label{chap:library:brauer}

Let $\ell$ be a prime, let $G$ be finite, let $k$ be an algebraically closed
field of characteristic $\ell$, and let $K$ be a field of characteristic
zero. Fix a multiplicative correspondence between the roots of unity in
$k$ and $K$ of orders dividing $|G|_{\ell'}$. We use it to lift the
eigenvalues of a modular representation and define its Brauer character.
For the usual complex valued construction, see, for example,
\cite[pp.~17--18]{Navarro1998}. The construction here uses only a finite
group of roots of unity in $K$, which need not be algebraically closed.

We prove additivity on short exact sequences, separation of irreducible
representations and linear independence of their Brauer characters.
These results will identify composition multiplicities from character
values in Chapter~\ref{chap:library:kzero}.

\section{Regular elements and roots of unity}
\label{lib:brauer:roots}

An element is $\ell$-regular if its order is prime to $\ell$. Write
$G_{\ell'}$ for these elements and
$\operatorname{CF}_{\ell'}(G,K)$ for the $K$-valued functions on
$G_{\ell'}$ invariant under conjugation by $G$. Set
\[
 m_G=|G|_{\ell'},\qquad
 \mu_{m_G}(k)=\{a\in k^\times:a^{m_G}=1\}.
\]
Here $|G|_{\ell'}$ is the largest divisor of $|G|$ prime to $\ell$.
It is a common exponent for the regular elements and need not be their
least common exponent.

\begin{definition}
\label{lib:brauer:root-correspondence}
A root correspondence for $G$ is a multiplicative isomorphism
$\iota_G:\mu_{m_G}(k)\to\mu_{m_G}(K)$.
\end{definition}

Since $k$ is algebraically closed and $\ell\nmid m_G$, its group of
$m_G$th roots is cyclic of order $m_G$. The existence of $\iota_G$
therefore requires $K$ to contain all $m_G$th roots of unity.
The correspondence $\iota_G$ is specified separately from the modular system in
Chapter~\ref{chap:library:reduction}. When restricting characters or
comparing character tables, the chosen correspondences must agree on the
roots involved.

\begin{implementation}
\module{ModularRep.BrauerRootConvention} defines
\lean{PrimeRegularRootEmbedding.AgreesOnRoots}: two correspondences agree
on every root whose order divides the exponent for the group being restricted
to. Agreement is reflexive and is transitive when the required exponents
divide one another. The theorem \lean{agreesOnRoots_of_commonRoot}
proves agreement for correspondences induced by one common root
identification. Separate choices of roots do not establish this agreement.
\end{implementation}

\begin{lemma}[Orders of eigenvalues at regular elements]
\label{lib:brauer:eigenvalues}
Let $\rho$ be a finite dimensional $k$-representation of $G$.
If $g\in G_{\ell'}$ and $a$ is a root of the characteristic polynomial of
$\rho(g)$, then $a^{m_G}=1$.
\end{lemma}
\begin{proof}
The order of $g$ divides $|G|$ and is prime to $\ell$, so it divides $m_G$.
Hence $g^{m_G}=1$. A root of the characteristic polynomial is an eigenvalue, so
there is $v\ne0$ with $\rho(g)v=av$. Applying the $m_G$th power gives
$v=a^{m_G}v$, and cancellation proves the assertion.
\end{proof}

\begin{implementation}
\lean{ModularRep.primeRegularExponent} is $m_G$.
\lean{ModularRep.PrimeRegularRootEmbedding} specifies a prime
\lean{p} and an equivalence between the groups defined by Mathlib's
\lean{rootsOfUnity}. The exponent identity and the
eigenvalue lemma are
\lean{ModularRep.PrimeRegularElement.pow_primeRegularExponent_eq_one} and
\lean{Representation.charpolyRoot_pow_primeRegularExponent_eq_one} in
\module{ModularRep.PrimeRegularEigenvalues}.
These two elementary results only need a field and a finite dimensional
representation. The field characteristic and algebraic closure enter the
subsequent construction of Brauer characters.
\end{implementation}

\section{The character of a modular representation}
\label{lib:brauer:construction}

For a finite dimensional $k$-representation $\rho$, let
$\mathcal R_\rho(g)$ be the multiset of roots of the characteristic
polynomial of $\rho(g)$, with algebraic multiplicities.

\begin{definition}
\label{lib:brauer:character-definition}
The Brauer character afforded by $\rho$, relative to $\iota_G$, is
\[
 \varphi_\rho(g)=\sum_{a\in\mathcal R_\rho(g)}\iota_G(a)
 \qquad(g\in G_{\ell'}).
\]
An irreducible Brauer character is one afforded by an irreducible finite
dimensional $k$-representation. Their set is denoted
$\IBr_{\iota_G}(G)$.
\end{definition}

The preceding lemma places every root in the domain of $\iota_G$.
Algebraic closure ensures that the characteristic polynomial splits over
$k$. The eigenvalues are lifted to $K$ before they are summed.
Their sum in $k$ is the trace $\tr_k(\rho(g))$.

\begin{proposition}[Invariance of Brauer characters]
\label{lib:brauer:invariance}
The function $\varphi_\rho$ is constant on conjugacy classes in
$G_{\ell'}$. Equivalent representations have equal Brauer characters.
For $\alpha\in\Aut(G)$,
\[
 \varphi_{\rho\circ\alpha}(g)=\varphi_\rho(\alpha(g)).
\]
\end{proposition}
\begin{proof}
Conjugating a group element conjugates its representing operator.
An equivalence of representations also conjugates the corresponding
operators. Conjugate operators have the same characteristic polynomial,
hence the same root multiset. The definition as a sum of lifted roots gives the first two
assertions. For the last assertion the operator on the left at $g$ is
exactly the operator on the right at $\alpha(g)$.
\end{proof}

\begin{implementation}
\module{ModularRep.BrauerCharacter} defines
\lean{Representation.brauerCharacterOfRootEmbedding} using Mathlib's
characteristic polynomials and their root multisets. It has values in
\lean{ModularRep.PrimeRegularClassFunction}. It uses a total function
\lean{PrimeRegularRootEmbedding.lift}, extended by zero outside the root
group, to map the root multiset. The theorem
\lean{Representation.lift_charpolyRootAsRootOfUnity} proves agreement with
the specified root equivalence on every root that occurs.
The invariance declarations are
\lean{Representation.charpoly_conj},
\lean{Representation.brauerCharacterOfRootEmbedding_iso} and
\lean{Representation.brauerCharacterOfRootEmbedding_twist}.
The conjugacy lemma specialises Mathlib's
\lean{LinearEquiv.charpoly_conj} to representation operators.
\lean{ModularRep.IBr} consists of class functions admitting an
irreducible \lean{FDRep} realisation. Equality of its elements is
pointwise equality of functions.
\end{implementation}

\begin{example}[Diagonal representations]
Suppose $g\in G_{\ell'}$ acts diagonally with entries
$a_1,\ldots,a_d\in\mu_{m_G}(k)$. Its Brauer value is
$\iota_G(a_1)+\cdots+\iota_G(a_d)$. If several entries coincide, every
occurrence contributes to the sum. At the identity all entries are one,
so $\varphi_\rho(1)=d$ in $K$. This remains the integer $d$ even when the
trace of the identity in $k$ vanishes because $\ell\mid d$.
\end{example}

\section{Additivity on exact sequences}
\label{lib:brauer:exactness}

\begin{lemma}[Characteristic polynomial on a subspace and quotient]
\label{lib:brauer:charpoly-factorisation}
Let $T$ be an endomorphism of a finite dimensional vector space $V$ over a
field, and let $U\leq V$ be $T$-invariant. If $\bar T$ is the induced
endomorphism of $V/U$, then
\[
 \det(XI-T)=\det(XI-T|_U)\det(XI-\bar T).
\]
\end{lemma}
\begin{proof}
Choose a basis of $U$, a basis of $V/U$ and lifts of the quotient basis.
Together they give a basis of $V$ in which $T$ has the form
\[
 \begin{pmatrix}A&B\\0&D\end{pmatrix}.
\]
The diagonal blocks represent $T|_U$ and $\bar T$. The determinant of the
resulting block upper triangular polynomial matrix is the product of the
two diagonal determinants.
\end{proof}

\begin{implementation}
\lean{LinearMap.charpoly_eq_restrict_mul_mapQ} in
\module{ModularRep.CharpolyInvariantSubmodule} constructs this basis and
identifies all four blocks. It uses Mathlib's basis construction
\lean{Module.Basis.sumQuot} and block matrix identity
\lean{Matrix.charpoly_fromBlocks_zero}\textsubscript{\texttt{21}}.
\end{implementation}

\begin{theorem}[Additivity on short exact sequences]
\label{lib:brauer:short-exact}
For a short exact sequence of finite dimensional $kG$-representations
\[
 0\longrightarrow U\overset{f}{\longrightarrow}V
   \overset{\pi}{\longrightarrow}W\longrightarrow0,
\]
the Brauer characters formed using $\iota_G$ satisfy
$\varphi_V=\varphi_U+\varphi_W$. In particular,
$\varphi_{U\oplus W}=\varphi_U+\varphi_W$.
\end{theorem}
\begin{proof}
First take an invariant subspace of $V$. By
Lemma~\ref{lib:brauer:charpoly-factorisation}, the root multiset of the
operator on $V$ is the sum of the root multisets on the subspace and its
quotient. Applying the same lift to each occurrence and summing proves
additivity. This step uses no additive property of $\iota_G$.

For a short exact sequence, $\im f=\ker \pi$ is invariant. Injectivity of
$f$ identifies $U$ with $\im f$, and surjectivity of $\pi$ identifies
$V/\im f$ with $W$. Both identifications intertwine the actions. The
first part and Proposition~\ref{lib:brauer:invariance} give the assertion.
The case of a direct sum follows in the same way, or from the characteristic
polynomial of a block diagonal operator.
\end{proof}

\begin{implementation}
\module{ModularRep.BrauerCharacterExact} proves
\lean{Representation.brauerCharacterOfRootEmbedding_subrepresentation_quotient}
and \lean{Representation.brauerCharacterOfRootEmbedding_prod}.
The statement for the category of finite dimensional representations is
\lean{FDRep.brauerCharacterOfRootEmbedding_shortExact}.
\lean{FDRep.shortExact_underlying_functions} gives injectivity,
surjectivity and exactness of the underlying linear maps. Intertwining
identities identify the subrepresentation and quotient with the first
and last terms of the sequence.
\end{implementation}

\begin{example}[An invariant filtration]
If $0=V_0\leq V_1\leq\cdots\leq V_r=V$ is a finite invariant filtration,
repeated exact additivity gives
$\varphi_V=\sum_{i=1}^r\varphi_{V_i/V_{i-1}}$.
For a representation given by block upper triangular matrices, the
diagonal blocks therefore suffice to calculate its Brauer character.
The blocks above the diagonal can distinguish a nonsplit extension from
a direct sum, even though their Brauer characters agree.
\end{example}

\section{Separation and linear independence}
\label{lib:brauer:traces}

The modular traces of nonisomorphic simple representations are linearly
independent. To pass from this fact to Brauer characters, we compare the
same sums of roots of unity in the two fields. Cyclotomic polynomials
provide the comparison. For linear independence of complex valued Brauer
characters, see \cite[Theorem~2.6]{Navarro1998}.

\begin{lemma}[Traces on regular parts]
\label{lib:brauer:regular-traces}
Let $\rho$ be a finite dimensional $k$-representation of $G$.
For every $g\in G$, write $g=g_\ell g_{\ell'}$ as commuting factors of
$\ell$-power order and order prime to $\ell$. Then
\[
 \tr_k(\rho(g))=\tr_k(\rho(g_{\ell'})).
\]
\end{lemma}
This trace identity also appears in \cite[Lemma~2.4]{Navarro1998}.
\begin{proof}
If $g_\ell$ has order dividing $\ell^a$, then in characteristic $\ell$
the operator $P=\rho(g_\ell)$ satisfies
$(P-I)^{\ell^a}=P^{\ell^a}-I=0$. Let $R=\rho(g_{\ell'})$.
The operators $P$ and $R$ commute. Hence $R(P-I)$ is nilpotent and has
trace zero. Since $RP=R+R(P-I)$, taking traces proves the claim.
\end{proof}

\begin{implementation}
\lean{Representation.character_eq_character_primeRegularPart} in
\module{ModularRep.ModularTraceRegularPart} uses the commuting primary
decomposition from \module{ModularRep.PrimeRegularPart} and Mathlib's trace
identity for a commuting nilpotent term.
\end{implementation}

\begin{proposition}[Brauer characters determine irreducible representations]
\label{lib:brauer:separation}
If two finite dimensional modular representations have equal Brauer
characters relative to $\iota_G$, then their $k$-valued trace functions
on $G$ are equal. If both representations are irreducible, they are
equivalent.
\end{proposition}
\begin{proof}
Choose a primitive $m_G$th root $\zeta\in k$. Its image
$\widehat\zeta=\iota_G(\zeta)$ is primitive in $K$. At a regular element,
write the two eigenvalue multisets as powers of $\zeta$. The difference
of the sums of these powers is the value of some polynomial
$F\in\Z[X]$ at $\zeta$. Equality of the Brauer values says
$F(\widehat\zeta)=0$. The minimal polynomial of $\widehat\zeta$ over
$\Q$ is the cyclotomic polynomial $\Phi_{m_G}$, so
$\Phi_{m_G}$ divides $F$ in $\Z[X]$. Since $m_G$ is prime to
$\ell$, the primitive root $\zeta$ is also a root of the corresponding
cyclotomic polynomial in $k$. Thus $F(\zeta)=0$.

The sum of the roots of the characteristic polynomial is the trace, so the traces
agree at regular elements. Lemma~\ref{lib:brauer:regular-traces} gives
agreement on all of $G$. The argument separating simple modules by their traces in
Proposition~\ref{lib:ordinary:separation}, applied over the algebraically
closed field $k$, proves the last assertion.
\end{proof}

\begin{implementation}
\module{ModularRep.CyclotomicRootSumReflection} proves that
equality of the lifted sums in $K$ implies equality of the sums in $k$.
\module{ModularRep.BrauerTraceRecovery} identifies these sums with the
root multisets used by the character definition.
\lean{ModularRep.FDRepSimpleClassKZero.brauerCharacterDeterminesModularTrace_of_rootEmbedding}
and
\lean{ModularRep.FDRepSimpleClassKZero.irreducibleBrauerCharacterInjectivity_of_rootEmbedding}
are the resulting theorems in \module{ModularRep.BrauerCharacterSeparation}.
The latter proves that simple modules with equal Brauer characters are
isomorphic. Trace recovery is proved from the root correspondence.
\end{implementation}

\begin{theorem}[Linear independence of irreducible Brauer characters]
\label{lib:brauer:independence}
The functions in $\IBr_{\iota_G}(G)$ are linearly independent over $K$
on $G_{\ell'}$.
\end{theorem}
\begin{proof}
Choose one representative of every simple $kG$-module class. There are
finitely many such classes: each is a composition factor of the finite
dimensional regular module. Denote the representatives by
$S_1,\ldots,S_s$ and their trace functions by $t_1,\ldots,t_s$.

First, the $t_i$ are linearly independent over $k$. For each $i$, Schur's
lemma and density on $\bigoplus_jS_j$ give an element of $kG$ acting as
an operator of trace one on $S_i$ and as zero on every other summand.
Evaluating any linear relation of trace functionals at this element
extracts its $i$th coefficient. Restriction of a linear functional on
$kG$ to the basis $G$ is injective, so the trace functions on $G$ are
independent too. Lemma~\ref{lib:brauer:regular-traces} then gives
independence of their restrictions to $G_{\ell'}$.

Independent functions admit a nonsingular square evaluation matrix.
Indeed, their evaluation vectors span $k^s$, since a linear functional
annihilating that span would give a relation between the functions.
Choose a basis of evaluation vectors, at regular elements
$x_1,\ldots,x_s$, and write
\[
 M_k=(t_i(x_j))_{i,j},\qquad \det M_k\ne0.
\]
Express each root multiset occurring in this matrix as powers of the
primitive root $\zeta$. Replacing each $\zeta$ by the class of $X$
defines a matrix $M$ over the common ring
$R=\Z[X]/(\Phi_{m_G})$. Evaluation at $\zeta$ gives $M_k$.
Evaluation at $\widehat\zeta$ gives the Brauer evaluation matrix
$M_K=(\varphi_{S_i}(x_j))_{i,j}$.

The map $R\to K$ is injective by the theorem on the minimal polynomial of a primitive root of unity. If $\det M_K=0$, then $\det M=0$ in $R$, and evaluating in
$k$ would give $\det M_k=0$. This contradiction shows
$\det M_K\ne0$. Evaluating a $K$-linear relation among the Brauer
characters at $x_1,\ldots,x_s$ now forces all coefficients to vanish.
Finally, separation and realisation identify this family indexed by
simple classes with the set of functions $\IBr_{\iota_G}(G)$.
\end{proof}

\begin{implementation}
\module{ModularRep.SimpleTraceLinearIndependence} proves independence of
the modular simple trace functions, and
\module{ModularRep.PrimeRegularEvaluationMatrix} obtains the regular
evaluation points. The coefficient argument is
\lean{ModularRep.CyclotomicDeterminantSpecialization.rootMultisetMatrix_det_ne_zero_transfer},
using \lean{CyclotomicOrder} and
\lean{cyclotomicSpecialization_injective} in the same namespace.
\lean{ModularRep.FDRepSimpleClassKZero.exists_simpleClass_primeRegular_brauer_det_ne_zero}
constructs the nonsingular Brauer evaluation matrix.
The final theorem is
\lean{ModularRep.FDRepSimpleClassKZero.irreducibleBrauerCharacters_linearIndependent}.
The determinant is first computed over the cyclotomic integer ring.
Its nonvanishing after evaluation in $K$ permits the final linear
algebra argument over $K$.
\end{implementation}

\begin{example}[Equality of composition multiplicities]
Suppose exact additivity gives
$\varphi_V=\sum_i a_i\varphi_i$ and
$\varphi_W=\sum_i b_i\varphi_i$, where $a_i,b_i$ are nonnegative integers
and the $\varphi_i$ are distinct irreducible Brauer characters.
If $\varphi_V=\varphi_W$, linear independence gives
$a_i=b_i$ for every $i$, using characteristic zero to recover equality
of integers from equality in $K$. This is the step used to prove
injectivity of the Brauer character map on the exact Grothendieck group.
It determines composition multiplicities without asserting that
$V\simeq W$.
\end{example}

\section{Compatible choices under pullback and restriction}
\label{lib:brauer:compatible-roots}

Let $f:H\to G$ be a homomorphism of finite groups with root
correspondences $\iota_H$ and $\iota_G$. A regular element of $H$ maps
to a regular element of $G$, so class functions can always be pulled back
by $f$. To compare this pullback with the character of $f^*\rho$, the
two lifts must agree on the roots of $\rho(f(h))$ for every
$h\in H_{\ell'}$.

\begin{proposition}[Pullback with compatible root correspondences]
\label{lib:brauer:compatible-pullback}
Let $f:H\to G$ be a homomorphism and let $\rho$ be a finite dimensional
$k$-representation of $G$. Suppose that $\iota_H(a)=\iota_G(a)$
for every eigenvalue $a$ of $\rho(f(h))$, for every $h\in H_{\ell'}$.
Then
\[
 \varphi_{f^*\rho}^{\,\iota_H}(h)
     =\varphi_\rho^{\,\iota_G}(f(h))
       \qquad(h\in H_{\ell'}).
\]
Such agreement holds if the two correspondences are restrictions of one
multiplicative equivalence $\mu_M(k)\simeq\mu_M(K)$, where both
$m_H$ and $m_G$ divide $M$.
\end{proposition}
\begin{proof}
The representing operator on both sides is $\rho(f(h))$. Its root
multiset is therefore the same on both sides, and the agreement of
the lifts gives equality term by term. A common root correspondence
gives that agreement on every root in the intersection of the two
root groups.
\end{proof}

\begin{implementation}
Compatibility on these eigenvalues is expressed by
\lean{Representation.BrauerRootLiftCompatibleAlong}.
\lean{Representation.brauerCharacterOfRootEmbedding_pullback_of_compatible}
proves the formula.
\module{ModularRep.BrauerCharacterCommonRootCompatibility} restricts one
common equivalence and proves the condition for any homomorphism.
Equality of the total functions $k\to K$ is a sufficient stronger
condition, but is unnecessary away from the occurring roots.
\end{implementation}

A group isomorphism $\phi:G\to H$ preserves $|G|_{\ell'}$ and therefore
transports a chosen root correspondence with the same function lifting
roots. Pullback along $\phi^{-1}$ then transports an irreducible Brauer
character together with an affording representation. The constructions
\lean{PrimeRegularRootEmbedding.alongMulEquiv} and
\lean{IrreducibleBrauerCharacter.alongMulEquiv} in
\module{ModularRep.BrauerCharacterEquivTransport} give this transport.
If an extension is also transported, the ambient and base group
isomorphisms must commute with the inclusion maps. A separately specified
target character must be identified with the transported base character.
These are the hypotheses of
\lean{Representation.Extension.BrauerCharacterExtensionWitness.alongMulEquivOfBase}
in \module{ModularRep.BrauerCharacterExtensionWitnessTransport}.

\begin{proposition}[Extending a root correspondence from a subgroup]
\label{lib:brauer:root-reselection}
Let $H\leq G$, and prescribe a root correspondence for $H$. If at least
one root correspondence for $G$ exists, one can choose a correspondence
for $G$ whose restriction agrees with the prescribed correspondence
on $\mu_{m_H}(k)$.
\end{proposition}
\begin{proof}
The divisibility $m_H\mid m_G$ permits restriction of any chosen ambient
correspondence to the smaller root group. If $j$ is this restriction, then
$j^{-1}\iota_H$ is an automorphism of the finite cyclic group
$\mu_{m_H}(k)$. An automorphism of a finite cyclic subgroup extends
to the ambient finite cyclic group: in power coordinates this is
surjectivity of reduction on the corresponding groups of units.
Compose the ambient correspondence with such an extension. Its
restriction is now the prescribed one.
\end{proof}

\begin{implementation}
\lean{ModularRep.PrimeRegularRootEmbedding.exists_ambient_agreeing_on_subgroup}
assumes existence of an ambient root correspondence, expressed by
\lean{Nonempty}, and uses
\lean{exists_mulAut_extension} in the same namespace. It adjusts an
ambient correspondence while retaining the prescribed subgroup correspondence.
\end{implementation}

There is a particularly direct choice for a normal $\ell$-subgroup
$N\trianglelefteq G$. Since $|G|=|G/N|\,|N|$, one has
$m_G=m_{G/N}$. A prescribed quotient correspondence can therefore be
transported to $G$ with the same root lifting map.
\lean{ModularRep.PrimeRegularRootEmbeddingPQuotient.exponent_quotient_eq},
\lean{ofPQuotient} and \lean{ofPQuotient_lift} implement this observation.

\section{Inflation and descent of irreducible Brauer characters}
\label{lib:brauer:descent}

For $\varphi\in\IBr_{\iota_G}(G)$, define $\ker\varphi=\ker\rho$,
where $\rho$ is any irreducible representation affording $\varphi$.
Proposition~\ref{lib:brauer:separation} makes this independent of the
choice of $\rho$, since equivalent representations have the same kernel.
Let $\pi:G\twoheadrightarrow H$ be a surjection.

\begin{theorem}[Inflation of irreducible Brauer characters]
\label{lib:brauer:quotient-equivalence}
Suppose that $\ell\nmid|\ker \pi|$, and that the root correspondences
are compatible along $\pi$ for every finite dimensional representation
of $H$. Inflation gives a bijection between
$\IBr_{\iota_H}(H)$ and
$\{\varphi\in\IBr_{\iota_G}(G):\ker \pi\leq\ker\varphi\}$.
If $\bar\varphi$ corresponds to $\varphi$, then
\[
 \varphi(g)=\bar\varphi(\pi(g))\qquad(g\in G_{\ell'}).
\]
\end{theorem}
\begin{proof}
Pullback along $\pi$ preserves irreducibility, and the pulled back
representation has $\ker \pi$ in its kernel. Root compatibility identifies
its Brauer character with the inflated function. Conversely, descend a
representation affording $\varphi$ on the same vector space, using
$\ker \pi\leq\ker\varphi$.
Representation descent preserves irreducibility, and compatible
pullback gives the asserted identity.

It remains to prove that pullback is injective on functions on $H_{\ell'}$. Every lift
$g$ of an $\ell$-regular $h$ is itself $\ell$-regular in this case.
Indeed, if $h$ has order $d$, then $g^d\in\ker \pi$, so the order of $g$ divides
$d\,|\ker \pi|$, which is prime to $\ell$. Hence $\pi$ is surjective on
regular elements and pullback of their class functions is injective.
The two character identities now show that inflation and descent are inverse.
\end{proof}

\begin{implementation}
\lean{ModularRep.IrreducibleBrauerCharacterSurjectiveDescent.KernelTrivialIBrAlong}
consists of irreducible Brauer characters admitting a realisation on
which $\ker \pi$ acts trivially. By character separation, this is the
condition $\ker \pi\leq\ker\varphi$. In the same namespace,
\lean{inflateToKernelTrivialIBrAlong}, \lean{descendIBrAlong} and
\lean{quotientIBrEquivKernelTrivialIBrAlong} construct the maps and the
bijection. The representation descent itself only requires surjectivity
and an irreducible realisation on which $\ker \pi$ acts trivially. Root
compatibility proves its character identity. The hypothesis that $\ker \pi$
has order prime to $\ell$ proves injectivity of pullback on class functions.
This hypothesis differs from the normal $\ell$-subgroup hypothesis in
the preceding construction of a root correspondence.
The classical inflation correspondence allows an arbitrary normal kernel
\cite[p.~39]{Navarro1998}. The theorem stated here uses the stronger
hypothesis that the kernel has order prime to $\ell$, as does the cited
Lean equivalence.
\end{implementation}

\section{Tensor actions and cyclic extensions}
\label{lib:brauer:actions-extensions}

For a modular linear character $\lambda:G\to k^\times$, define
$\widehat\lambda(g)=\iota_G(\lambda(g))$ on $G_{\ell'}$.
The value $\lambda(g)$ belongs to the required root group because
$g^{m_G}=1$. Multiplicativity of the root correspondence proves
$\widehat{\lambda\mu}=\widehat\lambda\widehat\mu$ and
$(\widehat\lambda)^\alpha=\widehat{\lambda\circ\alpha}$.

\begin{theorem}[Tensoring with a linear character]
\label{lib:brauer:tensor-source}
For every finite dimensional modular representation $V$ and every
$\lambda:G\to k^\times$, one has
\[
 \varphi_{V\otimes\lambda}=\widehat\lambda\,\varphi_V.
\]
\end{theorem}
\begin{proof}
At $g\in G_{\ell'}$, the representing operator on $V\otimes\lambda$
is $\lambda(g)\rho_V(g)$. Its eigenvalues, counted with multiplicity,
are $\lambda(g)a$ as $a$ runs through the eigenvalues of $\rho_V(g)$.
Multiplicativity of $\iota_G$ gives
$\iota_G(\lambda(g)a)=\widehat\lambda(g)\iota_G(a)$.
Summing proves the formula.
\end{proof}

This is the case of a one dimensional factor in the tensor product
formula for Brauer characters. See, for example,
\cite[proof of Theorem~2.23]{Navarro1998}.
The argument with invariant subspaces in
Proposition~\ref{lib:ordinary:linear-twist} also applies over $k$.
Thus $V\otimes\lambda$ is irreducible when $V$ is irreducible, and
$\widehat\lambda\varphi_V$ is then an irreducible Brauer character.

\begin{implementation}
The Lean construction uses
\lean{ModularRep.BrauerLinearTensorProductFormula} as a hypothesis.
\module{ModularRep.BrauerLinearCharacterTensorAction} defines
\lean{ModularRep.PrimeRegularRootEmbedding.liftedLinearCharacter} and proves
its multiplicative and automorphism identities. Given the tensor formula,
\lean{ModularRep.IrreducibleBrauerCharacter.linearTwist} constructs the
irreducible Brauer character. Its \lean{linearTwist_mul},
\lean{twist_linearTwist} and \lean{tensorFieldSemidirectAction} give the
action laws, with the inverse convention explained in
Section~\ref{lib:ordinary:actions}.
\end{implementation}

\begin{theorem}[Cyclic extension of a Brauer character]
Let $H$ be finite, let $N\trianglelefteq H$, and suppose that $H/N$
is cyclic. Choose root correspondences $\iota_N$ and $\iota_H$
compatible with restriction to $N$ for every finite dimensional
$k$-representation of $H$. If
$\varphi\in\IBr_{\iota_N}(N)$ is invariant under conjugation by $H$,
then there is $\widehat\varphi\in\IBr_{\iota_H}(H)$ with
$\widehat\varphi|_{N_{\ell'}}=\varphi$.
\end{theorem}
\begin{proof}
Choose an irreducible realisation $\rho$ of $\varphi$. Its conjugation
twists have the same Brauer character, so
Proposition~\ref{lib:brauer:separation} makes them equivalent to $\rho$.
The proof of Theorem~\ref{lib:ordinary:cyclic-principle} applies over
the algebraically closed field $k$ in any characteristic. It therefore
gives an extension $\widehat\rho$, which is irreducible because
$\rho$ is irreducible. Compatible roots identify the restriction of
its Brauer character with $\varphi$ by
Proposition~\ref{lib:brauer:compatible-pullback}.
\end{proof}

For the usual complex valued formulation and the representation extension
argument, see \cite[Theorem~8.12]{Navarro1998}.
There is no restriction on the order of the cyclic quotient relative to
$\ell$. The value field $K$ need not be algebraically closed, but the
root correspondences must satisfy the stated compatibility.

\begin{implementation}
The Lean construction assumes
\lean{Representation.BrauerCyclicExtensionPrinciple} over the
algebraically closed modular field $k$. It does not prove this
extension property from algebraic closure.
\module{ModularRep.BrauerCharacterExtensionBridge} proves
\lean{Representation.conjugationInvariant_of_brauerCharacter_fixed} and
\lean{Representation.exists_extension_realisation_of_ibr_fixed_cyclic_quotient}.
The latter gives an affording representation and a representation
extension. In \module{ModularRep.BrauerCharacterHomPullback},
\lean{Representation.Extension.brauerCharacterExtensionWitnessOfCompatible}
uses compatible roots to construct the extending irreducible Brauer character
and its restriction equality.
\end{implementation}
